% Retained preamble for this source, sliced MECHANICALLY from its own main TeX.
% Source file : full-original-native.tex (lines 1-122 of the file)
% The block below is verbatim author text; the wrapper class, document metadata and
% \begin{document} are omitted because the assembler supplies them.
% lines taken: 1, 2, 3, 4, 5, 6, 7, 8, 9, 10, 11, ... (122 lines total)
%~Mouliné par MaN_auto v.0.27.3 2022-08-08 16:38:20
\documentclass[AHL,Unicode,longabstracts,francais,published]{cedram}

\usepackage{float}
\usepackage{tikz}

\DeclareMathOperator{\Grass}{Grass}
\DeclareMathOperator{\covol}{covol}
\DeclareMathOperator{\GL}{GL}
\DeclareMathOperator{\GO}{O}
\DeclareMathOperator{\Supp}{Supp}
\DeclareMathOperator{\diag}{diag}
\DeclareMathOperator{\rg}{rg}
\DeclareMathOperator{\Stab}{Stab}



\usetikzlibrary{shapes.misc}
\tikzset{cross/.style={cross out, draw, minimum size=2*(#1-\pgflinewidth), inner sep=0pt, outer sep=0pt}}







\newcommand{\abs}[1]{\lvert#1\rvert}
\newcommand{\norm}[1]{\lVert#1\rVert}
\newcommand{\bracket}[1]{\langle#1\rangle}
\newcommand{\bcs}{\backslash}
\newcommand{\se}[2]{\{#1,\dots,#2\}}

\newcommand{\R}{\mathbb{R}}
\newcommand{\C}{\mathbb{C}}
\newcommand{\Z}{\mathbb{Z}}
\newcommand{\N}{\mathbb{N}}
\newcommand{\Q}{\mathbb{Q}}
\newcommand{\PP}{\mathbb{P}}
\newcommand{\QQ}{\overline{\mathbb{Q}}}
\newcommand{\NS}{\mathbb{S}}

\newcommand{\bw}{\mathbf{w}}
\newcommand{\be}{\mathbf{e}}
\newcommand{\bv}{\mathbf{v}}
\newcommand{\bp}{\mathbf{p}}

\newcommand{\cA}{\mathcal{A}}
\newcommand{\cM}{\mathcal{M}}
\newcommand{\cF}{\mathcal{F}}

\newcommand{\bbeta}{\widehat{\beta}}
\newcommand{\ggamma}{\widehat{\gamma}}

\newcommand{\eps}{\varepsilon}

\newcommand{\pgcd}{\mathrm{pgcd}}

\makeatletter
\def\editors#1{%
\def\editor@name{#1}
\if@francais
\def\editor@string{Recommand\'e par les \'editeurs \editor@name.}
\else
\def\editor@string{Recommended by Editors \editor@name.}
\fi}
\makeatother

%%%%%%%%%%%%%%%%%%%%%%%%%%%%%%%%%%%%%%%%%%%%%%%%%%%%%%%%%%

\graphicspath{{./figures/}}

\newcommand*{\mk}{\mkern -1mu}
\newcommand*{\Mk}{\mkern -2mu}
\newcommand*{\mK}{\mkern 1mu}
\newcommand*{\MK}{\mkern 2mu}

\hypersetup{urlcolor=purple, linkcolor=blue, citecolor=red}

\newcommand*{\romanenumi}{\renewcommand*{\theenumi}{\roman{enumi}}}
\newcommand*{\Romanenumi}{\renewcommand*{\theenumi}{\Roman{enumi}}}
\newcommand*{\alphenumi}{\renewcommand*{\theenumi}{\alph{enumi}}}
\newcommand*{\Alphenumi}{\renewcommand*{\theenumi}{\Alph{enumi}}}
\let\oldhat\hat
\renewcommand*{\hat}[1]{\mathchoice{\widehat{#1}}{\widehat{#1}}{\oldhat{#1}}{\oldhat{#1}}}
\let\oldexists\exists
\renewcommand*{\exists}{\mathrel{\oldexists}}
\let\oldforall\forall
\renewcommand*{\forall}{\mathrel{\oldforall}}
%%%%%%%%%%%%%%%%%%%%%%%%%%%%%%%%%%%%%%%%%%%%%%%%%%%%%%%

\title{Approximation diophantienne sur les quadriques}
\alttitle{Diophantine approximation on quadrics}

\subjclass{11J83, 11J87, 37A17}
\keywords{espaces de réseaux, points rationnels, groupes orthogonaux, théorème du sous-espace}


\author[\initial{N.} \lastname{de Saxcé}]{\firstname{Nicolas} \lastname{de Saxcé}}
\address{CNRS -- Université Sorbonne Paris Nord\\
LAGA / UMR 7539\\
99 avenue Jean-Baptiste Clément,\\
93430 Villetaneuse (France)}
\email{desaxce@math.univ-paris13.fr}


\begin{abstract}
Nous étudions l'exposant diophantien d'un point $x$ d'une hypersurface quadratique. Nous montrons notamment un analogue du théorème de Thue--Siegel--Roth, c'est-à-dire une formule pour l'exposant diophantien d'un point algébrique, et un analogue du résultat de Kleinbock et Margulis sur l'extrémalité des sous-variétés non dégénérées de l'espace affine.
\end{abstract}

\begin{altabstract}
We study the Diophantine exponent of a point $x$ on a quadric hypersurface. We show in particular an analogue of the Thue--Siegel--Roth theorem, that is to say a formula for the Diophantine exponent of an algebraic point, and an analogue of the result of Kleinbock and Margulis on the extremality of non-degenerate analytic manifolds in the of affine space.
\end{altabstract}


\datereceived{2021-07-09}
\daterevised{2022-03-30}
\dateaccepted{2022-05-12}

\editors{S. Gou\"ezel et X. Caruso}
%%%%%%%%%%%%%%%%%%%%%%%%%%%%%%%%%%%%%%%%%%%%%%%%%%%%

\dateposted{2022-11-14}
